% Generated draft: requires mathematical and source review.
\documentclass[11pt]{article}
\usepackage[margin=25mm]{geometry}
\usepackage{fontspec}
\setmainfont{Latin Modern Roman}
\usepackage{amsmath,amssymb,mathtools,unicode-math}
\setmathfont{Latin Modern Math}
\usepackage{enumitem,needspace,xurl,hyperref}
\setlength{\parindent}{0pt}
\setlength{\parskip}{5pt}
\setlength{\emergencystretch}{3em}
\newcommand{\N}{\mathbb{N}}
\newcommand{\Z}{\mathbb{Z}}
\newcommand{\Q}{\mathbb{Q}}
\newcommand{\R}{\mathbb{R}}
\newcommand{\C}{\mathbb{C}}
\newcommand{\F}{\mathbb{F}}
\newcommand{\A}{\mathbb{A}}
\newcommand{\PP}{\mathbb{P}}
\newcommand{\E}{\mathbb{E}}
\newcommand{\eps}{\varepsilon}
\newcommand{\dd}{\,\mathrm{d}}
\begin{document}
\section*{UnsolvedMath ID: 701303}
\textbf{Source problem code:} catalogue entry
\subsection{Definitions and mathematical statement}
For a finite group \(G\) with identity \(1_G\), interpret a nontrivial subgroup partition as a family \(H_1,\ldots,H_m\), with \(m\geq 2\), of nontrivial proper subgroups such that
\[
G=\bigcup_{i=1}^{m}H_i,
\qquad H_i\cap H_j=\{1_G\}\quad(i\ne j).
\]
Thus the sets \(H_i\setminus\{1_G\}\) partition \(G\setminus\{1_G\}\). This is the standard interpretation used here; the supplied statement does not explicitly define its partition convention.

Does there exist such a finite group \(G\) and such a partition for which
\[
|H_1|=\cdots=|H_m|
\quad\text{and}\quad
\exists i\in\{1,\ldots,m\}\ \exists x,y\in H_i:\ xy\ne yx?
\]
The last condition requires at least one partition subgroup to be nonabelian.

The supplied source states that any group satisfying these conditions must have prime exponent: there must exist a prime \(p\) such that every nonidentity element of \(G\) has order \(p\). This is retained as an attributed necessary condition, not as an independently verified result. The source also refers to Kourovka Notebook Problem 15.26, whose statement is not supplied.
\subsection{Short English statement}
Does a finite group admit a nontrivial partition into subgroups of equal order with at least one nonabelian subgroup? The partition is understood to cover every nonidentity element exactly once. The source reports that any example must have prime exponent and labels the existence question open.
\subsection{Sources}
\begin{itemize}
\item Supplied local UnsolvedMath registry record: permanent ID \(701303\), source code \texttt{catalogue entry}, title \emph{Kourovka Notebook Problem 17.1}. The supplied status is \texttt{open}; this label is not an independently verified literature assessment.
\item \emph{The Kourovka Notebook (Unsolved Problems in Group Theory)}, Problem 17.1, 17th issue (2010); the supplied background also identifies the 21st issue edition (2026), \texttt{arXiv:1401.0300v48}. The collection description names E. I. Khukhro and V. D. Mazurov as editors.
\item Attribution as supplied: A. Abdollahi is listed as proposer, with proposed year 2010; the statement itself begins with the attribution ``(I. M. Isaacs).'' Both attributions are retained without inferring their respective roles further.
\item I. M. Isaacs, \emph{Pacific J. Math.} \textbf{49} (1973), 109--116. The supplied statement cites this reference for the necessary prime-exponent condition; no article title is supplied.
\item Supplied comparison reference: Kourovka Notebook Problem 15.26. Its contents are not included in the record.
\end{itemize}
\end{document}
